\chapter{Cardiac Failure}

\section*{Definition and Overview}
Cardiac failure, also called heart failure, is a long-term condition in which the heart cannot pump blood strongly enough to meet the body's needs. It does not mean that the heart has stopped; rather, it is weakened, or too stiff to fill and empty effectively. Blood and fluid then build up in the lungs, legs, and other tissues, causing breathlessness, swelling, and fatigue. Heart failure can develop suddenly or gradually, is usually progressive, and requires lifelong management, although modern treatment can greatly improve symptoms, quality of life, and survival. It affects both the heart's ability to pump forward and the body's ability to cope with the resulting congestion.

\section*{Epidemiology}
Heart failure is common and becomes more so with age, affecting roughly 1 in 50 adults and about 1 in 10 people over the age of 75. It is one of the most frequent causes of hospital admission in older people, and readmission within months of discharge is common. Rates have been rising because populations are ageing and because better treatment of heart attacks means more people survive with weakened hearts. High blood pressure, coronary artery disease, and diabetes are the leading underlying causes, and the condition places a substantial burden on both patients and health services.

\section*{Aetiology and Pathophysiology}
Anything that damages the heart muscle, or chronically overloads it, can lead to heart failure:
\begin{itemize}
    \item \textbf{Coronary artery disease and heart attack:} the commonest cause; blocked arteries deprive heart muscle of oxygen, leaving scarred, weakened tissue
    \item \textbf{High blood pressure:} forces the heart to work harder for years, thickening and stiffening the muscle
    \item \textbf{Cardiomyopathy:} disease of the heart muscle itself, which may be inherited or caused by infections, alcohol, chemotherapy, or unknown factors
    \item \textbf{Heart valve disease:} narrowed or leaking valves overload the heart
    \item \textbf{Arrhythmias:} very fast or very slow heart rhythms reduce the efficiency of pumping
    \item \textbf{Other causes:} thyroid disease, kidney disease, anaemia, and inflammation of the heart muscle (myocarditis)
\end{itemize}
In heart failure, the heart's reduced pumping ability, or its stiffness on filling, means that blood is not moved forward properly; fluid accumulates in the lungs and body, and organs receive too little blood to function normally.

\section*{Clinical Features}
Symptoms may develop gradually or come on suddenly with an acute episode:
\begin{itemize}
    \item Breathlessness on exertion, on lying flat, or waking the person at night
    \item Swelling of the ankles, feet, and legs (oedema)
    \item Fatigue and weakness, even after light activity
    \item Rapid weight gain from fluid retention
    \item A persistent cough or wheeze, sometimes producing frothy or pink-tinged sputum
    \item Reduced appetite, bloating, and abdominal discomfort
    \item Confusion or reduced concentration in older people
\end{itemize}

\section*{Differential Diagnosis}
Many of the symptoms of heart failure also occur in other conditions, and misdiagnosis is common without testing:
\begin{itemize}
    \item Chronic lung disease, such as COPD or asthma, causing breathlessness and cough
    \item Pneumonia or other lung infections
    \item Pulmonary embolism, a blood clot in the lung
    \item Chronic kidney disease causing fluid retention and swelling
    \item Anaemia and thyroid disease causing fatigue and breathlessness
    \item Liver disease or venous problems causing leg swelling
    \item Obesity and poor physical fitness producing similar exercise intolerance
\end{itemize}

\section*{When to Seek Medical Care}
Anyone with new, worsening, or unexplained breathlessness, swelling, or fatigue should see a doctor. People with known heart failure should seek advice early if their weight rises rapidly or their symptoms worsen despite usual treatment.
\textbf{Call 000 (or local emergency services) for:}
\begin{itemize}
    \item Severe breathlessness at rest, or waking gasping for air
    \item Chest pain or pressure
    \item Fainting, collapse, or confusion
    \item Coughing up pink or bloody sputum
\end{itemize}

\section*{Diagnosis and Investigations}
Heart failure is diagnosed by combining the history, examination, and tests:
\begin{itemize}
    \item \textbf{History and examination:} asking about breathlessness, swelling, and risk factors; listening to the heart and lungs; and checking for fluid retention and elevated jugular venous pressure
    \item \textbf{Blood tests:} B-type natriuretic peptide (BNP or NT-proBNP), which rises when the heart is stretched, plus kidney, liver, and thyroid function and a full blood count
    \item \textbf{Electrocardiogram (ECG):} to detect previous heart attacks, arrhythmias, or strain patterns
    \item \textbf{Echocardiogram:} ultrasound showing the size of the heart, the strength of pumping (ejection fraction), and valve function
    \item \textbf{Chest X-ray:} to look for fluid in the lungs and enlargement of the heart
    \item \textbf{Further tests:} exercise testing, 24-hour heart rhythm monitoring, or cardiac MRI where the cause is unclear
\end{itemize}

\section*{Management}
Treatment aims to relieve symptoms, slow progression, and prevent hospital admissions:
\begin{itemize}
    \item \textbf{Diuretic medicines:} help the kidneys remove excess fluid, reducing swelling and breathlessness
    \item \textbf{Heart-protecting medicines:} drugs that relax blood vessels, slow the heart, and block harmful stress hormones; these reduce symptoms and improve survival
    \item \textbf{Lifestyle measures:} limiting salt, monitoring daily weight, staying as active as symptoms allow, and stopping smoking
    \item \textbf{Devices:} an implanted defibrillator or pacemaker for selected people with poor pumping function or dangerous rhythms
    \item \textbf{Procedures:} surgery or stenting for blocked arteries, valve repair, or, in severe cases, a mechanical heart pump or heart transplantation
    \item \textbf{Acute episodes:} hospital care with oxygen, intravenous medicines, and close monitoring
    \item \textbf{Cardiac rehabilitation and follow-up:} structured exercise, education, and regular specialist review
\end{itemize}

\section*{Complications}
Without good control, heart failure can lead to serious problems:
\begin{itemize}
    \item Fluid accumulation in the lungs causing severe breathlessness (pulmonary oedema)
    \item Arrhythmias, including dangerous rhythms that can cause sudden cardiac death
    \item Kidney damage from reduced blood flow and from some treatments
    \item Liver congestion and damage from back-pressure of fluid
    \item Blood clots, which can cause stroke or pulmonary embolism
    \item Muscle wasting, frailty, and poor nutrition over time
    \item Depression and reduced quality of life
\end{itemize}

\section*{Prognosis and Outcomes}
Heart failure is a serious condition, and its outlook depends on its cause, severity, and how well it is treated. With modern medication and structured follow-up, many people live for years with stable symptoms and a good quality of life. Some remain very well for long periods, while others have repeated hospital admissions. The condition is generally progressive, and in advanced cases a mechanical heart pump or transplantation may be considered. Prognosis is best when treatment is started early, medicines are taken reliably, and symptoms are monitored closely.

\section*{Prevention and Self-Care}
\begin{itemize}
    \item Take prescribed medicines exactly as directed, even when feeling well
    \item Weigh yourself daily and report a gain of more than about 2 kg over a few days
    \item Limit salt intake and reduce processed and salty foods
    \item Keep alcohol to a minimum, as it can weaken the heart further
    \item Stop smoking and keep blood pressure, cholesterol, and blood sugar controlled
    \item Stay gently active within your limits and join a cardiac rehabilitation programme if offered
    \item Have a plan for flare-ups, including who to call and when to seek urgent help
    \item Attend all follow-up appointments and report new symptoms promptly
\end{itemize}

\section*{Sources and Further Reading}
The following sources provide authoritative, up-to-date information on heart failure and its management.
\begin{itemize}
    \item NHS. \textit{Heart failure overview.} https://www.nhs.uk/conditions/heart-failure/
    \item Mayo Clinic. \textit{Heart failure: symptoms and causes.} https://www.mayoclinic.org/diseases-conditions/heart-failure/symptoms-causes/syc-20373142
    \item British Heart Foundation. \textit{Heart failure information.} https://www.bhf.org.uk/
    \item National Heart, Lung, and Blood Institute. \textit{Heart failure.} https://www.nhlbi.nih.gov/health/heart-failure
    \item MedlinePlus. \textit{Heart failure.} https://medlineplus.gov/heartfailure.html
\end{itemize}
